\begin{theoreme}
\label{IX.5.1}
Sea $g\colon S'\to S$ un morfismo de preesquemas que sea un morfismo de
descenso efectivo para la categoría fibrada de recubrimientos étales de
preesquemas (\cf \Ref{IX.4.9} y~\Ref{IX.4.12}). Supongamos que~$S$ es conexo
y que~$S'$, su cuadrado fibrado~$S''$ y su cubo fibrado~$S'''$ son sumas de
preesquemas conexos (como ocurre, por ejemplo, si~$S'$ es de tipo finito
sobre~$S$ localmente noetheriano y conexo). Elijamos, como antes, un punto
geométrico en cada componente conexa de~$S'$,~$S''$ y~$S'''$, ciertas clases
de caminos, un~$s_0'\in E'$ y, para cada~$s'\in E'$, un
\ifthenelse{\boolean{orig}}
{$s'\in E''$}
{$s''\in E''$}
cuyas dos imágenes en~$E'$ sean~$s_0'$ y~$s'$. ($E'$,~$E''$ y~$E'''$
designan, respectivamente, los conjuntos de componentes conexas de~$S'$,
~$S''$ y~$S'''$). Entonces el grupo fundamental de~$S$ en el punto geométrico
imagen de~$s_0'$ es canónicamente isomorfo al grupo de tipo galoisiano
engendrado por los~$\pi_{s'}=\pi_1(S',\underline{s}')$ ($s'\in E'$) y por
generadores~$g_{s''}$ ($s''\in E''$), sometidos a las relaciones
~\Ref{eq:IX.5.a},~\Ref{eq:IX.5.b} y~\Ref{eq:IX.5.c} anteriores, en las que
intervienen los elementos de los grupos
$\pi_{s''}=\pi_1(S'',\underline{s}'')$ y los elementos~$a_i^{s'''}$
($i=1,2,3$, $s'''\in E'''$) introducidos más arriba.
\end{theoreme}

\begin{corollaire}
\label{IX.5.2}
Supongamos que~$S'$ y~$S''$ no tienen más que un número finito de componentes
conexas y que los grupos fundamentales de las componentes conexas de~$S'$
son topológicamente finitamente generados. Entonces el grupo fundamental de
~$S$ es topológicamente finitamente generado.
\end{corollaire}

Así,
\marginpar{248}
demostraremos más adelante que el grupo fundamental de un esquema proyectivo
normal sobre un cuerpo algebraicamente cerrado es topológicamente finitamente
generado. Utilizando el lema de Chow y la normalización de los esquemas
algebraicos, se seguirá que el mismo resultado es cierto para todo esquema
propio sobre un cuerpo algebraicamente cerrado.

\begin{corollaire}
\label{IX.5.3}
Supongamos que~$S'$,~$S''$ y~$S'''$ no tienen más que un número finito de
componentes conexas, que los grupos fundamentales de las componentes conexas
de~$S'$ son topológicamente de \emph{presentación finita}, y que los grupos
fundamentales de las componentes conexas de~$S''$ son topológicamente
\emph{finitamente generados}. Entonces el grupo fundamental de~$S$ es
topológicamente de \emph{presentación finita}.
\end{corollaire}

Obsérvese que
\ifthenelse{\boolean{orig}}{\Ref{IX.4.9}}{\Ref{IX.4.10}}
puede expresarse, al restringirlo a los \emph{recubrimientos} étales,
diciendo que \emph{un morfismo finito, radicial y sobreyectivo de preesquemas
noetherianos induce un isomorfismo de los grupos fundamentales}; en términos
figurados, se puede decir, pues, que el grupo fundamental es un
\emph{invariante topológico} para los preesquemas.
Más generalmente, mediante~\Ref{IX.5.1} se puede explicitar el efecto sobre
el grupo fundamental de operaciones sobre los preesquemas, tales como el
«pinzamiento» de un preesquema a lo largo de un conjunto finito de puntos,
que tiene un significado topológico sencillo. Se obtiene, por ejemplo:

\begin{corollaire}
\label{IX.5.4}
Sean $g\colon S'\to S$ un morfismo finito de presentación finita y~$T$ un
subconjunto discreto de~$S$. Para cada~$s\in S$, sea~$n(s)$ el «número geométrico de
puntos» de la fibra~$g^{-1}(s)$ (que puede describirse también como el grado
separable de~$g^{-1}(s)$ sobre~$\kres(s)$, suma de los grados separables de sus
extensiones residuales). Supongamos que, para~$s\in S-T$, se tiene~$n(s)=1$. Para cada
~$s\in T$, sea~$K_s$ una extensión algebraicamente cerrada de~$\kres(s)$,
sea~$I_s$ el conjunto de los puntos geométricos de~$S'$ con valores en~$K_s$
(un conjunto con~$n(s)$ elementos), sea~$I_s'$ el complementario de un punto
elegido de~$I_s$ y, finalmente, sea~$I'$ la unión de los~$I_s'$. Supongamos
que~$S'$ es conexo. Entonces el grupo fundamental de~$S$ es isomorfo al grupo
de tipo galoisiano engendrado por el grupo fundamental de~$S'$ y por
generadores~$g_i$ ($i\in I'$), sin relaciones suplementarias.
\end{corollaire}

Los
\marginpar{249}
detalles de la demostración se dejan al lector; el enunciado obtenido no es
más que la traducción, al lenguaje de la teoría de grupos, del hecho de que la
categoría~$C$ de los recubrimientos étales de~$S$ es equivalente a la
categoría de los recubrimientos étales~$X'$ de~$S'$, \emph{provistos}, para
cada~$s\in T$, de un sistema transitivo de biyecciones entre las~$n(s)$ fibras
de~$X'$ en los puntos de~$g^{-1}(s)$ con valores en~$K_s$. (En esta forma
intrínseca, naturalmente, ya no es necesario suponer que~$S'$ es conexo).

\begin{exemple}
\label{IX.5.5}
Se demuestra fácilmente que la curva racional~$\PP^1_k$ sobre un cuerpo
algebraicamente cerrado~$k$ es simplemente conexa\footnote{\Cf
\Ref{XI.1.1}.}. Por tanto, el grupo fundamental de una curva racional
completa que tiene exactamente un punto doble con~$n$ ramas analíticas es el
grupo de tipo galoisiano libre engendrado por~$n-1$ generadores. Por ejemplo,
en el caso de un punto doble ordinario se obtiene el grupo fundamental
~$\hat\ZZ$, como se anunció en~(\Ref{I.11}~a)). En cambio, la existencia de
una cúspide (que es un punto «geométricamente unirramoso») no influye en el
grupo fundamental.
\end{exemple}

\begin{corollaire}
\label{IX.5.6}
Sea $g\colon S'\to S$ un morfismo de preesquemas universalmente submersivo,
con fibras geométricamente conexas, siendo~$S$ conexo. Entonces~$S'$ es
conexo y, eligiendo un punto geométrico~$s'$ de~$S'$ y denotando por~$s$ su
imagen en~$S$, el homomorfismo
$$
\pi_1(S',s')\to \pi_1(S,s)
$$
es \emph{sobreyectivo}. Si~$g$ es un morfismo de descenso efectivo para la
categoría fibrada de los recubrimientos étales de preesquemas
(\cf \emph{\Ref{IX.4.12}}), introduzcamos el punto geométrico
~$s^{\prime\prime}=\diag(s')$ de~$S^{\prime\prime}=S'\times_S S'$ y los dos
homomorfismos
$$
p_{1*},p_{2*}\colon \pi_1(S^{\prime\prime},s^{\prime\prime}) \to
\pi_1(S',s')
$$
inducidos por las dos proyecciones. Entonces~$\pi_1(S,s)$ es isomorfo al
conúcleo de este par de homomorfismos en la categoría de los grupos de tipo
galoisiano, es decir, al cociente de~$\pi_1(S',s')$ por el subgrupo normal
cerrado engendrado por los
\marginpar{250}
elementos de la forma
$p_{1*}(g^{\prime\prime})p_{2*}(g^{\prime\prime})^{-1}$, con
$g^{\prime\prime}\in \pi_1(S^{\prime\prime},s^{\prime\prime})$.
\end{corollaire}

En efecto, por~\Ref{IX.3.4} se sabe que el funtor
$X\mto X\times_SS'$ de los recubrimientos étales de~$S$ en los
recubrimientos étales de~$S'$ es plenamente fiel, lo que equivale a que el
homomorfismo de grupos fundamentales sea un epimorfismo (\Ref{V.6.9}). La
última afirmación es consecuencia inmediata de la descripción de
~\Ref{IX.5.1}.

\begin{remarque}
\label{IX.5.7}
No se sabe actualmente si el grupo fundamental de un esquema propio sobre un
cuerpo algebraicamente cerrado~$k$ es topológicamente de presentación
finita\footnote{Esto parece muy improbable en el caso de las curvas lisas de
género~$g\geq 2$ en característica~$p>0$. En cambio, cuando se sustituye
~$\pi_1$ por su mayor cociente primo con~$p$, parece que las técnicas bien
conocidas permiten dar una respuesta afirmativa, incluso sin suponer que el
esquema sea propio. \Cf un trabajo en preparación de J.P\ptbl Murre.}. Utilizando
~\Ref{IX.5.3}, una técnica bien conocida de secciones hiperplanas y la
desingularización de las superficies normales, se reduce al caso de una
\emph{superficie lisa sobre}~$k$. Esto permite, al menos, demostrar por
métodos trascendentes que la respuesta es afirmativa en característica~$0$
(sin necesidad de admitir la triangulabilidad de las variedades algebraicas
singulares). En característica~$p>0$, la dificultad principal parece residir
en el caso de las curvas, cuyo grupo fundamental solo se sabe que es un
cociente del que aparece en el caso clásico (\cf el exposé siguiente), pero
el núcleo por el que se toma el cociente se conoce muy mal.
\end{remarque}

\begin{remarque}
\label{IX.5.8}
Además de~\Ref{IX.5.4} y~\Ref{IX.5.6}, se podrían explicitar otros casos
particulares en los que~\Ref{IX.5.1} adopta una forma especialmente sencilla.
Un caso interesante es aquel en que~$S$ es el cociente de~$S'$ por un grupo
finito de automorfismos~$\Gamma$. \emph{Entonces la categoría de los
recubrimientos étales
\ifthenelse{\boolean{orig}}{de~$S'$}{de~$S$}
es equivalente a la categoría de los recubrimientos étales~$X'$ de~$S'$ en
los que el grupo~$\Gamma$ actúa de modo compatible con su acción sobre~$S'$,
de manera que para cada~$s'\in S'$ y cada~$g\in\Gamma_{s'}$} (donde
~$\Gamma_{s'}$ designa \emph{el grupo de inercia} de~$s'$ en~$\Gamma$),
\emph{$g$ actúa trivialmente sobre la fibra}~$X_{s'}^\prime$. Si~$S'$ es conexo,
este enunciado se interpreta del modo siguiente. Sea~$C'_0$ la categoría de
los recubrimientos étales de~$S'$ en los que~$\Gamma$ actúa de modo compatible
con su acción sobre~$S'$ (pero sin satisfacer necesariamente la condición
anterior sobre los grupos de inercia de los puntos de~$S'$). Se ve fácilmente
que es una categoría galoisiana (\Ref{V.5}) y que, para
\marginpar{251}
cada punto geométrico~$a'$ de~$S'$, el funtor fibra
$X'\mto X_{a'}'$ sobre~$C'_0$ es un funtor fundamental. Sea
~$\pi_1(S',\Gamma;a')=G$ el grupo de automorfismos de este funtor, provisto de
su topología habitual. Se tiene entonces una sucesión exacta canónica
$$
e\to \pi_1(S',a')\to G\to \Gamma\to e
$$
(caso particular de~\Ref{V.6.13}, donde se toma por~$S$ el recubrimiento
trivial~$S'\times\Gamma$ de~$S'$ definido por~$\Gamma$, sobre el cual se hace
actuar~$\Gamma$ de modo evidente). Además, para cada punto geométrico~$b'$ de
~$S'$, se tiene un isomorfismo
$\pi_1(S',\Gamma;b')\to G=\pi_1(S',\Gamma;a')$, definido salvo un
automorfismo interior procedente de~$\pi_1(S',a')$; y, como hay un
homomorfismo evidente de~$\Gamma_{b'}$ al primer miembro, se obtiene un homomorfismo
$$
u_{b'}\colon \Gamma_{b'}\to G \quoi,
$$
definido salvo un automorfismo interior (procedente
\ifthenelse{\boolean{orig}}
{de~$\pi_1(S',a')$,}
{de~$\pi_1(S',a')$),}
cuyo compuesto con el homomorfismo canónico~$G\to\Gamma$ es la inclusión
canónica~$\Gamma_{b'}\to\Gamma$. Hecho esto, \emph{el grupo fundamental
~$\pi_1(S,a)$ es canónicamente isomorfo al grupo cociente de
~$G=\pi_1(S',\Gamma;a')$ por el subgrupo normal cerrado engendrado por las
imágenes de los homomorfismos~$\Gamma_{b'}\to G$}. En particular, la imagen de
~$\pi_1(S',a')$ en~$\pi_1(S,a)$ es un subgrupo normal, y el cociente
correspondiente es isomorfo a un cociente de~$\Gamma$. Además, se puede
reducir el número de «relaciones» introducidas considerando, para cada
~$g\in\Gamma$, $g\neq e$, el subpreesquema~$S'_g$ de coincidencias de los
automorfismos
\ifthenelse{\boolean{orig}}{$\id_S$ y~$g$ de~$S$}{$\id_{S'}$ y~$g$ de~$S'$},
eligiendo un punto geométrico~$b'_{g,i}$ en cada componente conexa de~$S'_g$
y, después, uno de los homomorfismos correspondientes
~$\pi_1(S',\Gamma;b'_{g,i})\to G$, de donde se obtienen levantamientos
~$\bar{g}_i$ de~$g$ en~$G$. Basta entonces tomar el cociente de~$G$ por el
subgrupo normal cerrado de~$G$ engendrado por los~$\bar{g}_i$.

Cuando~$a'$ es fijo por~$\Gamma$, se ve fácilmente que~$\Gamma$ actúa de modo
natural sobre~$\pi_1(S',a')$ y que~$G$ se identifica con el producto
semidirecto correspondiente. Identificando entonces~$\Gamma$
\ifthenelse{\boolean{orig}}
{a}
{con}
un subgrupo de~$G$, se ve que entre las relaciones
\marginpar{252}
introducidas más arriba, al tomar~$b'=a'$, aparecen «$g=e$» para
~$g\in\ifthenelse{\boolean{orig}}{I}{\Gamma}$. Por tanto, \emph{si~$S'$ tiene
un punto geométrico~$a'$ fijo por}~$\ifthenelse{\boolean{orig}}{I}{\Gamma}$
(es decir, un punto~$s'$ cuyo grupo de inercia es~$\Gamma$), \emph{entonces
~$\pi_1(S,a)$ es un grupo cociente del grupo cociente de tipo galoisiano de
~$\pi_1(S',a')$ obtenido «haciendo triviales» las acciones de~$\Gamma$ sobre
~$\pi_1(S',a')$; y es incluso isomorfo a este último grupo si se supone que,
para cada~$g\in\ifthenelse{\boolean{orig}}{G}{\Gamma}$, el lugar de inercia
~$S_g'$ es conexo y, por tanto, pasa por el punto subyacente de~$a'$}.
Esta última afirmación está, en efecto, contenida en la segunda descripción
dada más arriba de las relaciones que se han de imponer a~$G$.

Este último resultado se aplica, en particular, si se toma por~$S'$ la
potencia cartesiana~$X^n$ de un preesquema conexo sobre un cuerpo
algebraicamente cerrado, y por~$\Gamma$ el grupo simétrico
~$\Gamma={\goth {S}}_n$, que actúa de la manera habitual, con lo que~$S$ es la
producto simétrico \ifthenelse{\boolean{orig}}{\nieme}{de grado~$n$}
\ifthenelse{\boolean{orig}}{de~$S$}{de~$X$}.
Tomando por~$a'$ un punto geométrico cuyo punto subyacente pertenece a la
diagonal, se cumplen las condiciones anteriores, pues todos los lugares de
inercia~$S_g'$ contienen la diagonal. Utilizando el hecho, demostrado en el
exposé siguiente, de que si~$X$ es propio y conexo sobre~$k$, el grupo
fundamental de~$X^n$ se identifica con~$\pi_1(X)^n$, se obtiene el siguiente
resultado curioso: \emph{si~$X$ es propio y conexo sobre un cuerpo
algebraicamente cerrado~$k$, el grupo fundamental de su producto simétrico
\ifthenelse{\boolean{orig}}{\nieme}{de grado~$n$}, $n\geq 2$, es isomorfo a la abelianización del grupo fundamental
de~$X$}. (Ignoro si se conoce el hecho análogo en topología algebraica; debería
poder establecerse mediante el mismo método de descenso). Tomemos, por
ejemplo, por~$X$ la curva racional~$X=\PP^1_k$; se obtiene una
\ifthenelse{\boolean{orig}}{\Nieme}{enésima} demostración de
que~$\PP^r_k$ es simplemente conexo, utilizando que~$\PP^1_k$ lo es. Tomemos
ahora por~$X$ una curva no singular sobre~$k$ y sea~$n\geq 2g-1$, de modo que
~$\Symm^n(X)$ está fibrado sobre la jacobiana~$J$, con fibras espacios
proyectivos, y tiene, por tanto (como se verá mediante los resultados de los
dos exposés siguientes), el mismo grupo fundamental que~$J$. Se recupera así,
sin una reducción por etapas, el conocido hecho de que \emph{el grupo fundamental de la
jacobiana de~$X$ es isomorfo a la abelianización del grupo fundamental de
~$X$}.
\end{remarque}
